\documentclass[11pt]{report} 

\usepackage{projectreport}
\setcounter{tocdepth}{3}
\newcommand{\name}{Eminalp Koyuncu - 2307882}
\newcommand{\course}{EE568 - Selected Topics on Electrical Machines}
\newcommand{\projecttitle}{Project-2: Motor Winding Design and Analysis}
\newcommand{\submissiondate}{May 2026}


% ------------------------------------------------------------------------------
% Document
% ------------------------------------------------------------------------------
\begin{document}

\maketitle

\tableofcontents
\newpage
% ------------------------------------------------------------------------------
% Top matter
% ------------------------------------------------------------------------------
\section*{Introduction}
\addcontentsline{toc}{section}{Introduction}


\newpage

\section*{Integral Slot Winding Design}
\addcontentsline{toc}{section}{Integral Slot Winding Design}

In this section, the winding design is conducted for two different integer-slot winding 72 slots, 6 pole and 3 phase machines, one for the full-pitched winding case and one for 11/12 short pitched case. 

\subsection*{Full-Pitched Winding}
In this section, the winding of 72 slot, 6 pole, 3 phase machine designed with full-pitched coils.

\addcontentsline{toc}{subsection}{Full-Pitched Winding}

For the given configuration, we first calculate the number of slots per pole per phase.

\begin{equation}
    Q = 72 \;\text{slots}
\end{equation}
\begin{equation}
    p = 6 \;\text{poles}
\end{equation}
\begin{equation}
    m = 3 \;\text{phases}
\end{equation}
\begin{equation}
    q = \frac{Q}{pm} = \frac{72}{6*3} = 4 \;\text{slots per pole per phase} 
\end{equation}

Also, we can find the electrical slot angle as,

\begin{equation}
    \alpha_u= \frac{360^\circ}{Q/(p/2)}=\frac{360^\circ}{72/3}=15^\circ
\end{equation}

So, each successive slot has a relative phase angle of $15^\circ$. If we denote the slot index by $i$, then, slot $i$ has a phase angle of $(i-1)*15^\circ$. Also, we know that each pole occupies 4 slots. If we place phase $A_{1+}$ to slot $i=1$ (0 degrees), then we should place phase $A_{1-}$ at an electrical phase of $180^\circ$, sice we have a full-pitch winding. Since the electrical slot angle is $15^\circ$, we should place the return line $A_{1-}$ at slot $i = (180/15)=13$. With the same logic, we can find the slot placement of each of the phases. The winding diagram in this case is given in \ref{tab:fullpitch}.


\begin{landscape}



\subsubsection*{Winding Diagram Layout}
\addcontentsline{toc}{subsubsection}{Winding Diagram Layout}
\thispagestyle{empty}

\vspace*{\fill}

\begin{center}
\resizebox{1.6\textwidth}{!}{
    \begin{tabular}{c|*{24}{c|}}
    Slot & 1&2&3&4&5&6&7&8&9&10&11&12&13&14&15&16&17&18&19&20&21&22&23&24 \\
    Angle& $0^\circ$&$15^\circ$&$30^\circ$&$45^\circ$&$60^\circ$&$75^\circ$&$90^\circ$&$105^\circ$&$120^\circ$&$135^\circ$&$150^\circ$&$165^\circ$&$180^\circ$&$195^\circ$&$210^\circ$&$225^\circ$&$240^\circ$&$255^\circ$&$270^\circ$&$285^\circ$&$300^\circ$&$315^\circ$&$330^\circ$&$345^\circ$\\

    &
    \cellcolor{red!20}$A_{1+}$&
    \cellcolor{red!20}$A_{2+}$&
    \cellcolor{red!20}$A_{3+}$&
    \cellcolor{red!20}$A_{4+}$&
    \cellcolor{blue!20}$C_{5-}$&
    \cellcolor{blue!20}$C_{6-}$&
    \cellcolor{blue!20}$C_{7-}$&
    \cellcolor{blue!20}$C_{8-}$&
    \cellcolor{green!20}$B_{1+}$&
    \cellcolor{green!20}$B_{2+}$&
    \cellcolor{green!20}$B_{3+}$&
    \cellcolor{green!20}$B_{4+}$&
    \cellcolor{red!20}$A_{5-}$&
    \cellcolor{red!20}$A_{6-}$&
    \cellcolor{red!20}$A_{7-}$&
    \cellcolor{red!20}$A_{8-}$&
    \cellcolor{blue!20}$C_{1+}$&
    \cellcolor{blue!20}$C_{2+}$&
    \cellcolor{blue!20}$C_{3+}$&
    \cellcolor{blue!20}$C_{4+}$&
    \cellcolor{green!20}$B_{5-}$&
    \cellcolor{green!20}$B_{6-}$&
    \cellcolor{green!20}$B_{7-}$&
    \cellcolor{green!20}$B_{8-}$ \\

    &
    \cellcolor{red!20}$A_{5+}$&
    \cellcolor{red!20}$A_{6+}$&
    \cellcolor{red!20}$A_{7+}$&
    \cellcolor{red!20}$A_{8+}$&
    \cellcolor{blue!20}$C_{1-}$&
    \cellcolor{blue!20}$C_{2-}$&
    \cellcolor{blue!20}$C_{3-}$&
    \cellcolor{blue!20}$C_{4-}$&
    \cellcolor{green!20}$B_{5+}$&
    \cellcolor{green!20}$B_{6+}$&
    \cellcolor{green!20}$B_{7+}$&
    \cellcolor{green!20}$B_{8+}$&
    \cellcolor{red!20}$A_{1-}$&
    \cellcolor{red!20}$A_{2-}$&
    \cellcolor{red!20}$A_{3-}$&
    \cellcolor{red!20}$A_{4-}$&
    \cellcolor{blue!20}$C_{5+}$&
    \cellcolor{blue!20}$C_{6+}$&
    \cellcolor{blue!20}$C_{7+}$&
    \cellcolor{blue!20}$C_{8+}$&
    \cellcolor{green!20}$B_{1-}$&
    \cellcolor{green!20}$B_{2-}$&
    \cellcolor{green!20}$B_{3-}$&
    \cellcolor{green!20}$B_{4-}$ \\
    \end{tabular}
}

\captionof{table}{Winding Diagram for the Full-Pitched Integer Slot Winding Design}
\label{tab:fullpitch}
\end{center}

\fillandplacepagenumber

\end{landscape}

\subsubsection*{Distribution Factor, Pitch Factor and Winding Factor}
\addcontentsline{toc}{subsubsection}{Distribution Factor, Pitch Factor and Winding Factor}

The distribution factor for $n'th$ harmonic can be calculated as follows.
\begin{equation}
k_{d}(n)=\frac{\sin{(qn\frac{\alpha_u}{2})}}{q\sin{(n\frac{\alpha_u}{2})}}
\end{equation}
Also, the pitch factor for $n'th$ harmonic can be calculated as follows. 

\begin{equation}
    k_p(n)=\sin{(n\frac{\lambda}{2})}
\end{equation}

In these equations, $n$ is the harmonic number and $\lambda$ is the coil pitch. In the full pitch case, we have $k_p = 1$ for all the harmonics since $\lambda=180^\circ$. The distribution factor for the fundamental component, $3^{rd}$ and $5^{th}$ components can be calculated as follows.

\begin{equation}
k_{d}(1)=\frac{\sin{(4*\frac{20^\circ}{2})}}{4\sin{(\frac{20^\circ}{2})}} = 0.925
\end{equation}

\begin{equation}
k_{d}(3)=\frac{\sin{(4*3*\frac{20^\circ}{2})}}{4*\sin{(3\frac{\alpha_u}{2})}} = 0.433
\end{equation}

\begin{equation}
k_{d}(5)=\frac{\sin{(4*5*\frac{20^\circ}{2})}}{4\sin{(5*\frac{20^\circ}{2})}} = -0.112
\end{equation}

Thus, the winding factor $k_w = k_d*k_p$ for the fundamental, $3^{rd}$ and $5^{th}$ components can be calculated as follows.

\begin{equation}
    k_w(1) = 0.925
\end{equation}

\begin{equation}
    k_w(3) = 0.433 
\end{equation}

\begin{equation}
    k_w(5) = -0.112 
\end{equation}




\subsubsection*{Comments}
\addcontentsline{toc}{subsubsection}{Comments}

From the winding factor, the distributed coil causes to drop the effect of $3^{rd}$ harmonic to 43.3\% of its magnitude and $5^{th}$ harmonic to 11.2\% with a phase difference of $180^\circ$ due to the minus sign. The tradeoff is that the fundamental component magnitude drops to its 92.5\% due to the effect of the distribution. In this case, we have no effect of the coil pitch on the winding factor since the coil is full-pitch. The winding factor can be better (lower for the harmonics) with the introduction of the short-pitch winding. 


\subsection*{11/12 Short-Pitched Winding}
\addcontentsline{toc}{subsection}{11/12 Short-Pitched Winding}

For the short-pitch case, the slot angle and the number of slots per pole per phase is the same as the full-pitch case. The only difference is the coil pitch, which can be calculated as,

\begin{equation}
    \lambda=\frac{11}{12}*180^\circ=165^\circ
\end{equation}

So, the return line of each coil is phase shifted $165^\circ$ electrically. We know the slot angle is $15^\circ$, so, if we place phase $A_{1+}$ to slot $i=1$, then we should place the return line $A_{1-}$ to slot $i=(165/15) +1=12$. With a similar logic for other phases, we obtain the winding diagram given in \ref{tab:shortpitch}.



\begin{landscape}

\subsubsection*{Winding Diagram Layout}
\addcontentsline{toc}{subsubsection}{Winding Diagram Layout}
\thispagestyle{empty}

\vspace*{\fill}

\begin{center}
\resizebox{1.6\textwidth}{!}{
    \begin{tabular}{c|*{24}{c|}}
    Slot & 1&2&3&4&5&6&7&8&9&10&11&12&13&14&15&16&17&18&19&20&21&22&23&24 \\
    Angle& $0^\circ$&$15^\circ$&$30^\circ$&$45^\circ$&$60^\circ$&$75^\circ$&$90^\circ$&$105^\circ$&$120^\circ$&$135^\circ$&$150^\circ$&$165^\circ$&$180^\circ$&$195^\circ$&$210^\circ$&$225^\circ$&$240^\circ$&$255^\circ$&$270^\circ$&$285^\circ$&$300^\circ$&$315^\circ$&$330^\circ$&$345^\circ$\\

    &
    \cellcolor{red!20}$A_{1+}$&
    \cellcolor{red!20}$A_{2+}$&
    \cellcolor{red!20}$A_{3+}$&
    \cellcolor{red!20}$A_{4+}$&
    \cellcolor{blue!20}$C_{5-}$&
    \cellcolor{blue!20}$C_{6-}$&
    \cellcolor{blue!20}$C_{7-}$&
    \cellcolor{blue!20}$C_{8-}$&
    \cellcolor{green!20}$B_{1+}$&
    \cellcolor{green!20}$B_{2+}$&
    \cellcolor{green!20}$B_{3+}$&
    \cellcolor{green!20}$B_{4+}$&
    \cellcolor{red!20}$A_{5-}$&
    \cellcolor{red!20}$A_{6-}$&
    \cellcolor{red!20}$A_{7-}$&
    \cellcolor{red!20}$A_{8-}$&
    \cellcolor{blue!20}$C_{1+}$&
    \cellcolor{blue!20}$C_{2+}$&
    \cellcolor{blue!20}$C_{3+}$&
    \cellcolor{blue!20}$C_{4+}$&
    \cellcolor{green!20}$B_{5-}$&
    \cellcolor{green!20}$B_{6-}$&
    \cellcolor{green!20}$B_{7-}$&
    \cellcolor{green!20}$B_{8-}$ \\

    &
    \cellcolor{red!20}$A_{6+}$&
    \cellcolor{red!20}$A_{7+}$&
    \cellcolor{red!20}$A_{8+}$&
    \cellcolor{blue!20}$C_{1-}$&
    \cellcolor{blue!20}$C_{2-}$&
    \cellcolor{blue!20}$C_{3-}$&
    \cellcolor{blue!20}$C_{4-}$&
    \cellcolor{green!20}$B_{5+}$&
    \cellcolor{green!20}$B_{6+}$&
    \cellcolor{green!20}$B_{7+}$&
    \cellcolor{green!20}$B_{8+}$&
    \cellcolor{red!20}$A_{1-}$&
    \cellcolor{red!20}$A_{2-}$&
    \cellcolor{red!20}$A_{3-}$&
    \cellcolor{red!20}$A_{4-}$&
    \cellcolor{blue!20}$C_{5+}$&
    \cellcolor{blue!20}$C_{6+}$&
    \cellcolor{blue!20}$C_{7+}$&
    \cellcolor{blue!20}$C_{8+}$&
    \cellcolor{green!20}$B_{1-}$&
    \cellcolor{green!20}$B_{2-}$&
    \cellcolor{green!20}$B_{3-}$&
    \cellcolor{green!20}$B_{4-}$& 
    \cellcolor{red!20}$A_{8+}$ \\
    \end{tabular}
}

\captionof{table}{Winding Diagram for the 11/12 Short-Pitched Integer Slot Winding Design}
\label{tab:shortpitch}
\end{center}

\fillandplacepagenumber

\end{landscape}



\subsubsection*{Distribution Factor, Pitch Factor and Winding Factor}

\addcontentsline{toc}{subsubsection}{Distribution Factor, Pitch Factor and Winding Factor}

In this case, the distribution factor is the same as the full-pitch case since the number of slots per pole per phase is the same. 

\begin{equation}
k_{d}(1)=\frac{\sin{(4*\frac{20^\circ}{2})}}{4\sin{(\frac{20^\circ}{2})}} = 0.925
\end{equation}

\begin{equation}
k_{d}(3)=\frac{\sin{(4*3*\frac{20^\circ}{2})}}{4*\sin{(3\frac{\alpha_u}{2})}} = 0.433
\end{equation}

\begin{equation}
k_{d}(5)=\frac{\sin{(4*5*\frac{20^\circ}{2})}}{4\sin{(5*\frac{20^\circ}{2})}} = -0.112
\end{equation}

The difference comes from the pitch factor.


\begin{equation}
    k_p(1)=\sin{(\frac{165^\circ}{2})}=0.991
\end{equation}


\begin{equation}
    k_p(3)=\sin{(3\frac{165^\circ}{2})}=-0.924
\end{equation}


\begin{equation}
    k_p(5)=\sin{(5\frac{165^\circ}{2})}=0.793
\end{equation}

Hence, the winding factors for the fundamental, $3^{rd}$ and $5^{th}$ components become,

\begin{equation}
    k_w(1) = 0.925*0.991=0.917
\end{equation}

\begin{equation}
    k_w(3) = 0.433 *(-0.924)=-0.400
\end{equation}

\begin{equation}
    k_w(5) = (-0.112)*0.793= -0.089
\end{equation}



\subsubsection*{Comments}
\addcontentsline{toc}{subsubsection}{Comments}

In this case, utilizing the short-pitch winding, we obtain a better performance in terms of the harmonics rejection, with a very small loss in the fundamental component. Also, with the short pitch winding, a phase difference of $180^\circ$ is added on the third harmonic. 



\section*{Fractional-Slot Winding Design for Reference Design 2}
\addcontentsline{toc}{section}{Fractional-Slot Winding Design for Reference Design 2}

\subsection*{Phase Angle of Induced Voltage in Each Slot}
\addcontentsline{toc}{subsection}{Phase Angle of Induced Voltage in Each Slot}

\subsection*{Phasor Diagram for One Phase and Winding Factor}
\addcontentsline{toc}{subsection}{Phasor Diagram for One Phase and Winding Factor}

\subsection*{Complete Winding Diagram}
\addcontentsline{toc}{subsection}{Complete Winding Diagram}

\begin{landscape}
\thispagestyle{empty}

\vspace*{\fill}

\begin{center}
\resizebox{1.6\textwidth}{!}{
    \begin{tabular}{c|*{15}{c|}}
    Slot & 1&2&3&4&5&6&7&8&9&10&11&12&13&14&15 \\
    Angle& $0^\circ$&$120^\circ$&$240^\circ$&$0^\circ$&$120^\circ$&$240^\circ$&$0^\circ$&$120^\circ$&$240^\circ$&$0^\circ$&$120^\circ$&$240^\circ$&$0^\circ$&$120^\circ$&$240^\circ$\\
    &
    \begin{tabular}{@{}c|c@{}}
        \cellcolor{green!20}$B_{5-}$ & \cellcolor{red!20}$A_{1+}$ 
    \end{tabular}
    &
    \begin{tabular}{@{}c|c@{}}
        \cellcolor{red!20}$A_{1-}$ & \cellcolor{blue!20}$C_{1+}$
    \end{tabular}
    &
    \begin{tabular}{@{}c|c@{}}
        \cellcolor{blue!20}$C_{1-}$ & \cellcolor{green!20}$B_{1+}$
    \end{tabular}
    &
    \begin{tabular}{@{}c|c@{}}
        \cellcolor{green!20}$B_{1-}$ & \cellcolor{red!20}$A_{2+}$
    \end{tabular}
    &
    \begin{tabular}{@{}c|c@{}}
        \cellcolor{red!20}$A_{2-}$ & \cellcolor{blue!20}$C_{2+}$
    \end{tabular}
    &
    \begin{tabular}{@{}c|c@{}}
        \cellcolor{blue!20}$C_{2-}$ & \cellcolor{green!20}$B_{2+}$
    \end{tabular}
    &
    \begin{tabular}{@{}c|c@{}}
        \cellcolor{green!20}$B_{2-}$ & \cellcolor{red!20}$A_{3+}$
    \end{tabular}
    &
    \begin{tabular}{@{}c|c@{}}
        \cellcolor{red!20}$A_{3-}$ & \cellcolor{blue!20}$C_{3+}$
    \end{tabular}
    &
    \begin{tabular}{@{}c|c@{}}
        \cellcolor{blue!20}$C_{3-}$ & \cellcolor{green!20}$B_{3+}$
    \end{tabular}
    &
    \begin{tabular}{@{}c|c@{}}
        \cellcolor{green!20}$B_{3-}$ & \cellcolor{red!20}$A_{4+}$
    \end{tabular}
    &
    \begin{tabular}{@{}c|c@{}}
        \cellcolor{red!20}$A_{4-}$ & \cellcolor{blue!20}$C_{4+}$
    \end{tabular}
    &
    \begin{tabular}{@{}c|c@{}}
        \cellcolor{blue!20}$C_{4-}$ & \cellcolor{green!20}$B_{4+}$
    \end{tabular}
    &
    \begin{tabular}{@{}c|c@{}}
        \cellcolor{green!20}$B_{4-}$ & \cellcolor{red!20}$A_{5+}$
    \end{tabular}
    &
    \begin{tabular}{@{}c|c@{}}
        \cellcolor{red!20}$A_{5-}$ & \cellcolor{blue!20}$C_{5+}$
    \end{tabular}
    &
    \begin{tabular}{@{}c|c@{}}
        \cellcolor{blue!20}$C_{5-}$ & \cellcolor{green!20}$B_{5+}$
\end{tabular}
\end{tabular}
}

\captionof{table}{Winding Diagram for the Fractional-Slot Winding Design}
\label{tab:placeholder}
\end{center}

\fillandplacepagenumber

\end{landscape}


\section*{Analytical Modelling}
\addcontentsline{toc}{section}{Analytical Modelling}

\subsection*{Electrical Loading}
\addcontentsline{toc}{subsection}{Electrical Loading}

\subsection*{Carter's Coefficient and Effective Air-Gap Clearance}
\addcontentsline{toc}{subsection}{Carter's Coefficient and Effective Air-Gap Clearance}

\subsection*{Effective Axial Length}
\addcontentsline{toc}{subsection}{Effective Axial Length}

\subsection*{Air-Gap Peak Magnetic Flux Density}
\addcontentsline{toc}{subsection}{Air-Gap Peak Magnetic Flux Density}

\subsection*{Magnetic Loading, Specific Machine Constant, Shear Stress and Expected Torque}
\addcontentsline{toc}{subsection}{Magnetic Loading, Specific Machine Constant, Shear Stress and Expected Torque}

\subsection*{Flux Per Pole}
\addcontentsline{toc}{subsection}{Flux Per Pole}


\section*{2D FEA Modelling}
\addcontentsline{toc}{section}{2D FEA Modelling}

\subsection*{Geometry}
\addcontentsline{toc}{subsection}{Geometry}

\subsection*{Meshing}
\addcontentsline{toc}{subsection}{Meshing}

\subsection*{Boundary Conditions}
\addcontentsline{toc}{subsection}{Boundary Conditions}

\subsection*{Solver Settings}
\addcontentsline{toc}{subsection}{Solver Settings}


\subsection*{Results}
\addcontentsline{toc}{subsection}{Results}

\subsubsection*{Flux Density Distribution}
\addcontentsline{toc}{subsubsection}{Flux Density Distribution}

\subsubsection*{Air-Gap Flux Density Distribution}
\addcontentsline{toc}{subsubsection}{Air-Gap Flux Density Distribution}

\subsubsection*{Leakage Flux Paths}
\addcontentsline{toc}{subsubsection}{Leakage Flux Paths}


\subsubsection*{Flux-Per-Pole}
\addcontentsline{toc}{subsubsection}{Flux-Per-Pole}


\subsubsection*{Bonus: Cogging Torque}
\addcontentsline{toc}{subsubsection}{Bonus: Cogging Torque}




\end{document}